Táto práca sa zaoberá návrhom a optimalizáciou konvolučných neurónových sietí (\acrshort{cnn}) pre odhad intenzity slnečného žiarenia 
na základe obrazových dát oblohy. Hlavným cieľom je analyzovať vplyv rozmerov vstupných obrázkov na presnosť regresného 
modelu a identifikovať optimálnu konfiguráciu vstupných dát pre rôzne architektúry vlastných CNN modelov.

V rámci práce boli navrhnuté tri architektúry konvolučných neurónových sietí s rôznou komplexnosťou.
 Pre každú architektúru boli systematicky testované rôzne rozmery vstupných obrázkov spolu s kombináciami hyperparametrov, 
 ako sú rýchlosť učenia, veľkosť dávky, dropout a počet neurónov v plne prepojenej vrstve.

Modely boli trénované a vyhodnocované pomocou metrík strednej absolútnej chyby (\acrshort{mae}), strednej 
kvadratickej chyby (\acrshort{mse}) a 
koreňovej strednej kvadratickej chyby (\acrshort{rmse}). Súčasne bol analyzovaný čas trénovania a počet potrebných epoch, čím sa 
zohľadnila aj výpočtová náročnosť jednotlivých konfigurácií.

Výsledky ukazujú, že vhodná voľba rozmeru vstupných obrázkov má významný vplyv na presnosť modelu aj jeho výpočtovú 
efektívnosť. Na základe experimentov boli identifikované optimálne rozmery vstupov pre jednotlivé architektúry CNN, 
ktoré dosahujú najnižšiu chybu odhadu slnečného žiarenia.